% TOP_PROBLEM: {"id": 30007171, "title": "Exact maximal limit-cycle count for quintic classical Lienard systems", "category_id": 11, "category": {"id": 11, "name": "dynamical_systems", "display_name": "Dynamical Systems", "description": "Problems about long-term behavior of deterministic systems, Hamiltonian dynamics, and periodic orbits.", "slug": "dynamical-systems", "order_index": 11, "created_at": "2026-07-31T15:26:25.671Z"}, "status": "open"}
% FIRST_POSTED: 2026-10-06
% ATTEMPT_STATUS: solved
% OPENAI_CLAIM: {"schema_version":1,"source_repo":"https://github.com/openai/math","source_commit":"adc7f1241b42e322a6451854ab7e4b4c146bf78a","release_date":"2026-10-06","match":"full","resolution":"proved","independently_reviewed":false,"second_pass":"agreed","families":[{"family":"143","title":"Hilbert's sixteenth problem: uniform bounds for limit cycles","subject":"Dynamical systems and ergodic theory","manuscripts":[{"title":"Two limit cycles for quintic Liénard systems","pdf_path":"preprints/two-limit-cycles-for-quintic-lienard-systems-September-24-2026/two-limit-cycles-for-quintic-lienard-systems-September-24-2026.pdf","readme_path":"preprints/two-limit-cycles-for-quintic-lienard-systems-September-24-2026/README.md","date":"2026-09-24","theorem_ref":"thm:main (sections/introduction.tex:13)","statement_sources":[{"path":"preprints/two-limit-cycles-for-quintic-lienard-systems-September-24-2026/build/sections/introduction.tex","line":13,"label":"thm:main"}]}],"lean":{"doc_path":"lean/docs/143.md","comparators":["lean/ComparatorChallenges/QuinticLienard.json","lean/ComparatorChallenges/QuinticLienard.lean"],"declarations":["OAI.QuinticLienard.main"],"covers_main_theorem":true},"reasoning_trace":null}]}
\documentclass[11pt]{article}
\usepackage[margin=25mm]{geometry}
\usepackage{fontspec}
\setmainfont{Latin Modern Roman}
\usepackage{amsmath,amssymb,mathtools}
\usepackage{unicode-math}
\setmathfont{Latin Modern Math}
\usepackage{enumitem}
\usepackage{xurl,hyperref,bookmark}
\hypersetup{colorlinks=true,urlcolor=blue}
\setcounter{secnumdepth}{0}
\setlength{\parindent}{0pt}
\setlength{\parskip}{5pt}
\setlength{\emergencystretch}{3em}
\begin{document}
\section*{Exact maximal limit-cycle count for quintic classical Lienard systems}

\subsection{OpenAI's claim}
A classical Liénard system with arbitrary real primitive polynomial of degree at most five is claimed to have at most two isolated periodic orbits, and two are claimed to occur. The cited manuscript argues for that result; its relation to this problem is specified below.

\subsection{Mathematical statement of the claimed result}
\paragraph{Mathematical claim 1.} For any real polynomial $F$ of degree at most five, consider
\[
\dot x=y-F(x),\qquad\dot y=-x.
\]
OpenAI claims at most two limit cycles in $\mathbb R^2$, with examples attaining two. A cycle means the geometric image of a nonconstant periodic trajectory isolated among periodic orbits; multiplicity and stability do not affect the count. Periodic trajectories of a center are excluded.

\subsection{Outline of the argument}
Manuscript 1 (family 143, 2026-09-24): A quadratic midpoint model tracks fitted slope and curvature by transport equations. Endpoint variation, Riccati linearization, and a Schwarzian identity control the width dependence of the model coefficients, allowing the two-cycle argument without the odd-polynomial restriction of the older theorem.

\subsection{Scope relative to this problem}
Manuscript 1 (family 143, 2026-09-24): The problem selects the mathematical statement summarized here, including its stated hypotheses and parameter range. The cited theorem gives a claimed proof of that assertion in full.

\subsection{Formal verification}
For family 143, the scope page documents the following selections. The selection concerns degree-at-most-five Liénard systems: two isolated cycles are the maximum and that maximum occurs. It is not the arbitrary-degree planar polynomial bound. The documented statement covers the main claim of manuscript 1. Comparator files: QuinticLienard.lean. Declarations: OAI.QuinticLienard.main.

This site has not executed the Lean files or checked their dependencies. These descriptions report the scope of the linked declarations, rather than a new verification of the proofs.

\subsection{Status caveat}
OpenAI attributes these results to an internal model. Its release notice cautions that results lacking full formalization can contain mistakes. This site has not independently assessed the proofs; a solved label records the scope of the reported claim.

\subsection{Sources}
\begin{itemize}
\item OpenAI, \emph{Two limit cycles for quintic Liénard systems}, 2026-09-24; citation key \texttt{OAI:two-limit-cycles-for-quintic-lienard-systems-September-24-2026}. \href{https://github.com/openai/math/blob/adc7f1241b42e322a6451854ab7e4b4c146bf78a/preprints/two-limit-cycles-for-quintic-lienard-systems-September-24-2026/two-limit-cycles-for-quintic-lienard-systems-September-24-2026.pdf}{PDF (GitHub)}; \href{https://github.com/openai/math/raw/adc7f1241b42e322a6451854ab7e4b4c146bf78a/preprints/two-limit-cycles-for-quintic-lienard-systems-September-24-2026/two-limit-cycles-for-quintic-lienard-systems-September-24-2026.pdf}{PDF (direct)}; \href{https://github.com/openai/math/blob/adc7f1241b42e322a6451854ab7e4b4c146bf78a/preprints/two-limit-cycles-for-quintic-lienard-systems-September-24-2026/README.md}{Manuscript page and citation}; \href{https://github.com/openai/math/blob/main/preprints/two-limit-cycles-for-quintic-lienard-systems-September-24-2026/two-limit-cycles-for-quintic-lienard-systems-September-24-2026.pdf}{Latest version}; Statement source: \href{https://github.com/openai/math/blob/adc7f1241b42e322a6451854ab7e4b4c146bf78a/preprints/two-limit-cycles-for-quintic-lienard-systems-September-24-2026/build/sections/introduction.tex\#L13}{thm:main}.
\item Family 143: \href{https://github.com/openai/math/blob/adc7f1241b42e322a6451854ab7e4b4c146bf78a/lean/docs/143.md}{Lean scope}.
\item \href{https://github.com/openai/math/blob/adc7f1241b42e322a6451854ab7e4b4c146bf78a/lean/ComparatorChallenges/QuinticLienard.json}{Comparator QuinticLienard.json}.
\item \href{https://github.com/openai/math/blob/adc7f1241b42e322a6451854ab7e4b4c146bf78a/lean/ComparatorChallenges/QuinticLienard.lean}{Comparator QuinticLienard.lean}.
\item \href{https://github.com/openai/math/blob/adc7f1241b42e322a6451854ab7e4b4c146bf78a/CONTENTS.md}{OpenAI catalogue entry}.
\end{itemize}
\end{document}
